% ID: SEL-D0035-OA-E02
% ESTADO: cotejada. Segunda pasada visual y borrador compilado revisados.
% Fuente: Selectividad/matematicas_aplicadas/titular_septiembre_EXAMEN 1-ANDALUCÍA.pdf
% SHA-256: 7a09a02641ea8de345b94f49b27515e58436e032c7401867f86dc523490a3ea5
% Página: 1; opción A; ejercicio original 2.

Sea la función $f(x)=x^3+ax^2+bx$.
\begin{enumerate}[label=\alph*)]
\item \textbf{(1 punto)} Halle $a$ y $b$ sabiendo que la función tiene un mínimo en el punto de abscisa $x=-1$ y un punto de inflexión en el punto de abscisa $x=-2$.
\item \textbf{(1.5 puntos)} Para $a=6$ y $b=9$, halle los puntos de corte con los ejes, estudie la monotonía y extremos y esboce la gráfica de la función.
\end{enumerate}
